\documentclass[10pt,a4paper]{amsart}
\usepackage[margin=19mm]{geometry}
\usepackage{amsmath,amssymb,mathtools}
\usepackage[scr=rsfs,cal=euler]{mathalfa}
\usepackage{xfrac}
\usepackage[unicode,hidelinks,bookmarks=false]{hyperref}
\newcommand{\ie}{i.e.\ }
\newcommand{\cf}{cf.\ }
\newcommand{\resp}{respectively\ }
\newcommand{\Subsection}{\subsection}
\providecommand{\nobreakdash}{\nobreak\hbox{-}}
\setlength{\emergencystretch}{2em}
\multlinegap0pt
\usepackage{bm}
\usepackage{bbm}
\usepackage{dsfont}
\usepackage{xspace}
\usepackage{cancel}
\usepackage{mathtools}
\usepackage{amsthm}
\makeatletter
\@ifundefined{thm}{\newtheorem{thm}{Theorem}}{}
\@ifundefined{lem}{\newtheorem{lem}[thm]{Lemma}}{}
\makeatother
\newcommand{\Ad}{\mathrm{Ad}}
\newcommand{\ad}{\mathrm{ad}}
\newcommand{\fr}[1]{\mathfrak{#1}}
\newcommand{\largewedge}{\bigwedge}
\newcommand{\ox}{\otimes}
\newcommand{\pr}{\prime}
\newcommand{\rT}{\mathrm T}
\newcommand{\sm}[1]{\mathsf{#1}}
\title[Poisson-Lie Group structures on semidirect products]{Poisson-Lie Group structures on semidirect products}
\author{Floris Elzinga, Makoto Yamashita}
\date{Source: arXiv:2203.14552v2 (CC BY 4.0)}
\begin{document}
\maketitle

\noindent\emph{Source and licence.} Poisson-Lie Group structures on semidirect products. Version arXiv:2203.14552v2 (CC BY 4.0), distributed under the Creative Commons Attribution 4.0 International licence, as recorded in the arXiv licence metadata for this exact version and shown on the abstract page https://arxiv.org/abs/2203.14552v2. Full rights notice: licenses/arxiv-2203.14552-CC-BY-4.0.txt.\par\smallskip

\noindent\textit{Editorial scope.} Counted unit: the Thm whose source label is \texttt{thm:PLGS} in the article (source0.tex lines 362--364, source label \texttt{thm:PLGS}), delivered together with the article's own complete original proof of it (source0.tex lines 465--543, one \texttt{proof} environment of 79 lines). No mathematics is written, repaired or re-derived: statement and proof are the article's own text line for line. The proof is detached from the statement in the source: the article states the result at lines 362--364 and prints its proof at lines 465--543, and the proof environment's own header names the statement, so the association is the article's own. Required context retained uncounted: eq:ADEAZ (lines 319--325); eq:ADEA (lines 321--324); eq:LGPB (lines 346--351); eq:psi-i-y-i-invariance (lines 374--377); lem:Pi-corresp-to-eta (lines 386--389); eq:AEFLP (lines 424--429). Every definition, display and label of those lines is retained unchanged. Omitted from this package and not redistributed: 1--318, 325--345, 352--361, 365--373, 378--385, 390--423, 430--464, 544--1804. Nothing inside a retained excerpt is omitted or altered: every retained body is delivered exactly as the article's lines are written, so no omission receipt of any kind accompanies this package. Citations: the retained excerpts carry no citation command at all, so no bibliography entry of the article is transcribed and no reference is redistributed. Reference adaptation: none was needed and none was made; every reference inside the delivered unit resolves inside it, so no unresolved reference or citation is produced. Printed slips kept literal: no printed word, symbol, hypothesis or number of the counted statement or of its proof was altered. Layout and class: the article's own class is \texttt{amsart} with packages including fontenc, inputenc, lmodern, microtype, fullpage, amssymb, physics, tensor, stmaryrd, hyperref, amsrefs; the wrapper is the lane's pinned \texttt{amsart} editorial wrapper at 10pt on A4, which restates the theorem-like environments the retained text needs and adds the article's own macro definitions that the retained text needs (\texttt{\textbackslash Ad}, \texttt{\textbackslash ad}, \texttt{\textbackslash fr}, \texttt{\textbackslash largewedge}, \texttt{\textbackslash ox}, \texttt{\textbackslash pr}, \texttt{\textbackslash rT}, \texttt{\textbackslash sm}), with extra wrapper packages (bm, bbm, dsfont, xspace, cancel, mathtools). The delivered numbering is the unit's own. Assets: no retained span contains a figure environment, a graphics-inclusion command, a PDF-page inclusion, a table or a TikZ picture; no image file is copied into this package, the package declares no asset dependency and no image file is redistributed with it. Licence: version arXiv:2203.14552v2 of this article is distributed under the Creative Commons Attribution 4.0 International licence, as recorded in the arXiv licence metadata for this exact version and shown on the abstract page https://arxiv.org/abs/2203.14552v2. A full rights notice is delivered at licenses/arxiv-2203.14552-CC-BY-4.0.txt. Characters: every retained excerpt is pure ASCII, so the delivered engine, XeLaTeX with its default Latin Modern text font, reports no missing character.
\begin{lem}
The action of $g = (v, b)$ for $\Ad^{E}$ is as follows,
\begin{align}
\label{eq:ADEAZ} \Ad^{E}_g (\psi) &= \Ad^\vee_b \psi && (\psi \in \fr{b}^0),\\
 \label{eq:ADEA} \Ad^{E}_g (y) &= \Ad_b y - \ad^\vee ( \Ad_b y ) (v) && (y \in \fr{b}).
\end{align}
\end{lem}

\begin{align}
\label{eq:LGPB}
\left\{ \widetilde{X_1} , \widetilde{X_2} \right\} &= \widetilde{\left[ X_1 , X_2 \right]},&
\left\{ \widetilde{X} , \pi^*(f_1) \right\} &= \pi^*\left( \sm{a}(X) f \right),&
\left\{ \pi^*(f_1) , \pi^*(f_2) \right\} &= 0.
\end{align}

\begin{equation}
\label{eq:psi-i-y-i-invariance}
\sum_i \Ad^\vee_b \psi^i \ox P_{\fr{c}}\Ad_b y_i = \sum_i \psi^i \ox y_i \quad (b \in B).
\end{equation}

\begin{lem}
\label{lem:Pi-corresp-to-eta}
The Poisson bivector $\Pi \in \Gamma(E, \largewedge^2 \rT E)$ for \eqref{eq:LGPB} is given by $\Pi_g = (\rT R_g)_e^{\ox 2} \eta(g)$.
\end{lem}

\begin{lem}
With $g \in E$ presented by $(v, b) \in \fr{b}^0 \times B$, we have
\begin{equation}
\eta_0(g) = \frac{1}{2} \left\langle v , \left[ y_i , y_j \right] \right\rangle \psi^i \wedge \psi^j - \Ad^\vee_b \psi^i \wedge \ad^\vee ( P_{\fr{b}}\Ad_b y_i ) (v) . \label{eq:AEFLP}
\end{equation}
\end{lem}

\begin{thm}\label{thm:PLGS}
The Poisson bracket on $E$ characterized by \eqref{eq:LGPB}, together with the semidirect product group structure, defines a Poisson--Lie group structure.
\end{thm}

\begin{proof}[Proof of Theorem \ref{thm:PLGS}]
By Lemma \ref{lem:Pi-corresp-to-eta}, it is enough to show that $\eta$ satisfies the cocycle condition
\[
\eta(g h) = \eta(g) + \left( \Ad^{E}_g \ox \Ad^{E}_g \right) \eta(h),
\]
or equivalently,
\[
\left( \Ad^{E}_g \ox \Ad^{E}_g \right) \eta(h) = \eta(g h) - \eta(g) .
\]
We proceed by computing the left-hand side for the two pieces $\eta_0$ and $\eta_{\fr{b}}$ separately, and then compare.
In the following we write $g = (v, b)$ and $h = (w, b^\pr)$.

First we expand $\left( \Ad^{E}_g \ox \Ad^{E}_g \right) \eta_0(h)$ as
\begin{multline*}
\frac{1}{2} \left\langle w , \Ad_{b^\pr} \mathopen{}\left[ y_i , y_j \right] \right\rangle \Ad^\vee_{b b^\pr} \psi^i \wedge \Ad^\vee_{b b^\pr} \psi^j 
= \frac{1}{2} \left\langle \Ad^\vee_b w , \Ad_{b b^\pr} \mathopen{} \left[ y_i , y_j \right] \right\rangle \Ad^\vee_{b b^\pr} \psi^i \wedge \Ad^\vee_{b b^\pr} \psi^j \\
= \frac{1}{2} \left\langle v + \Ad^\vee_b w , \Ad_{b b^\pr} \mathopen{} \left[ y_i , y_j \right] \right\rangle \Ad^\vee_{b b^\pr} \psi^i \wedge \Ad^\vee_{b b^\pr} \psi^j - \frac{1}{2} \left\langle v , \Ad_{b b^\pr} \mathopen{} \left[ y_i , y_j \right] \right\rangle \Ad^\vee_{b b^\pr} \psi^i \wedge \Ad^\vee_{b b^\pr} \psi^j .
\end{multline*}
We recognize the first term as $\eta_0(g h)$, so we focus on the second term.
If we apply our formula \eqref{eq:AEFLP} to it, we find
\[
\frac{1}{2} \left\langle v , \Ad_{b b^\pr} \mathopen{} \left[ y_i , y_j \right] \right\rangle \Ad^\vee_{b b^\pr} \psi^i \wedge \Ad^\vee_{b b^\pr} \psi^j = \frac{1}{2} \left\langle v , \left[ y_i , y_j \right] \right\rangle \psi^i \wedge \psi^j - \Ad_{b b^\pr}^* \psi^i \wedge \ad^\vee ( P_{\fr{b}}\Ad_{b b^\pr} y_i ) (v).
\]
The relation \eqref{eq:AEFLP} also implies that
\[
\frac{1}{2} \left\langle v , \left[ y_i , y_j \right] \right\rangle \psi^i \wedge \psi^j = \frac{1}{2} \left\langle v , \Ad_b \mathopen{} \left[ y_i , y_j \right] \right\rangle \Ad^\vee_b \psi^i \wedge \Ad^\vee_b \psi^j + \Ad^\vee_b \psi^i \wedge \ad^\vee ( P_{\fr{b}} \Ad_b y_i ) (v) .
\]
Combining the two yields
\begin{multline*}
\frac{1}{2} \left\langle v , \Ad_{b b^\pr} \mathopen{} \left[ y_i , y_j \right] \right\rangle \Ad^\vee_{b b^\pr} \psi^i \wedge \Ad^\vee_{b b^\pr} \psi^j =\\
 \frac{1}{2} \left\langle v , \Ad_b \mathopen{} \left[ y_i , y_j \right] \right\rangle \Ad^\vee_b \psi^i \wedge \Ad^\vee_b \psi^j + \Ad^\vee_b \psi^i \wedge \ad^\vee ( P_{\fr{b}} \Ad_b y_i ) (v)\\
  - \Ad^\vee_{b b^\pr} \psi^i \wedge \ad^\vee ( P_{\fr{b}} \Ad_{b b^\pr} y_i ) (v) .
\end{multline*}
The first term on the right-hand side is precisely $\eta_0(g)$, so we now have
\begin{multline}\label{eq:g-on-eta-0-h}
\left( \Ad^{E}_g \ox \Ad^{E}_g \right) \eta_0(h) =\\
 \eta_0(g h) - \eta_0(g) - \Ad^\vee_b \psi^i \wedge \ad^\vee ( P_{\fr{b}} \Ad_b y_i ) (v) + \Ad^\vee_{b b^\pr} \psi^i \wedge \ad^\vee ( P_{\fr{b}} \Ad_{b b^\pr} y_i ) (v) .
\end{multline}

Second, we see that $\left( \Ad^{E}_g \ox \Ad^{E}_g \right) \eta_{\fr{b}}(h) $ is equal to
\[
\Ad^\vee_{b b^\pr} \psi^i \wedge \left[ \Ad_b P_{\fr{b}} \Ad_{b^\pr} y_i - \ad^\vee ( \Ad_b P_{\fr{b}} \Ad_{b^\pr} y_i ) (v) \right]
\]
by \eqref{eq:ADEAZ} and \eqref{eq:ADEA}.
To obtain projections in the right places, write
\[
\Ad_b P_{\fr{b}} \Ad_{b^\pr} y_i = P_{\fr{b}} \Ad_b P_{\fr{b}} \Ad_{b^\pr} y_i = P_{\fr{b}} \Ad_b \mathopen{} \left( 1 - P_{\fr{c}} \right) \Ad_{b^\pr} y_i .
\]
We proceed by exploiting the duality again,
\[
\begin{split}
\Ad^\vee_{b b^\pr} \psi^i \wedge \left[ P_{\fr{b}} \Ad_{b b^\pr} y_i - P_{\fr{b}} \Ad_b P_{\fr{c}} \Ad_{b^\pr} y_i \right] &= \Ad^\vee_{b b^\pr} \psi^i \wedge P_{\fr{b}} \Ad_{b b^\pr} y_i - \Ad^\vee_b \psi^i \wedge P_{\fr{b}} \Ad_b y_i \\
&= \eta_{\fr{b}}(g h) - \eta_{\fr{b}}(g) .
\end{split}
\]
This implies
\[
\left( \Ad^{E}_g \ox \Ad^{E}_g \right) \eta_{\fr{b}}(h) = \eta_{\fr{b}}(g h) - \eta_{\fr{b}}(g) - \Ad^\vee_{b b^\pr} \psi^i \wedge \ad^\vee ( P_{\fr{b}} \Ad_{b b^\pr} y_i - P_{\fr{b}} \Ad_b P_{\fr{c}} \Ad_{b^\pr} y_i ) (v),
\]
and by invariance \eqref{eq:psi-i-y-i-invariance} we have
\begin{equation}\label{eq:g-on-eta-b-h}
\begin{split}
\left( \Ad^{E}_g \ox \Ad^{E}_g \right) \eta_{\fr{b}}(h)
&= \eta_{\fr{b}}(g h) - \eta_{\fr{b}}(g) \\
&\quad - \Ad^\vee_{b b^\pr} \psi^i \wedge \ad^\vee ( P_{\fr{b}} \Ad_{b b^\pr} y_i ) (v) + \Ad^\vee_b \psi^i \wedge \ad^\vee ( P_{\fr{b}} \Ad_b y_i ) (v) .
\end{split}
\end{equation}

Combining \eqref{eq:g-on-eta-0-h} and \eqref{eq:g-on-eta-b-h}, we find
\[
\begin{split}
\left( \Ad^{E}_g \ox \Ad^{E}_g \right) \left( \eta_0(h) + \eta_{\fr{b}}(h) \right) &= \eta_0(g h) - \eta_0(g) + \eta_{\fr{b}}(g h) - \eta_{\fr{b}}(g)\\
&\quad  + \Ad^\vee_{b b^\pr} \psi^i \wedge \ad^\vee ( P_{\fr{b}} \Ad_{b b^\pr} y_i ) (v)- \Ad^\vee_{b b^\pr} \psi^i \wedge \ad^\vee ( P_{\fr{b}} \Ad_{b b^\pr} y_i ) (v) \\
&\quad- \Ad^\vee_b \psi^i \wedge \ad^\vee ( P_{\fr{b}} \Ad_b y_i ) (v) + \Ad^\vee_b \psi^i \wedge \ad^\vee ( P_{\fr{b}} \Ad_b y_i ) (v) \\
&= \eta(g h) - \eta(g) ,
\end{split}
\]
as desired.
\end{proof}

\end{document}
